% ============ 引言(第 0 章, 原书 pp.1--18 之前半) ============
% 定理沿用原书编号 0.1, 0.4, 0.5; 例 0.6, 0.7.
\chapter*{引言}\index{yinyan@引言(introduction)}\index{leiyulun@类域论(class field theory)}

% 引言内定理按 0.x 编号
\renewcommand{\thetheorem}{0.\arabic{theorem}}
\setcounter{theorem}{0}

L'objet de la th\'eorie du corps de classes est de montrer comment les extensions ab\'eliennes d'un corps de nombres alg\'ebriques $K$ peuvent \^etre d\'etermin\'ees par des \'el\'ements tir\'es de la connaissance de $K$ lui-m\^eme; ou, si l'on veut pr\'esenter les choses en termes dialectiques, comment un corps poss\`ede en soi les \'el\'ements de son propre d\'epassement.\footnote{类域论的目标是展示一个代数数域 $K$ 的 Abel 扩张如何能够由取自 $K$ 自身知识的元素来确定; 或者, 如果愿意用辩证的措辞来说, 一个域如何在自身中就包含着超越自身的那些元素. \cite{Chevalley1940}.}

类域论的目标, 是用域本身的算术来描述局部域或整体域的 Galois 扩张. 对 Abel 扩张, 该理论在大约 1850 至 1930 年间由 Kronecker, Weber, Hilbert, Takagi, Artin, Hasse 等人发展起来. 对非 Abel 扩张, 理论应有之形态, 最早见于 1967 年 Langlands 致 Weil 的一封信\cite{LanglandsWeil1967}. 近些年来, 非 Abel 的局部域与函数域情形已取得许多进展, 但数域情形进展较少. 约自 1980 年起, Abel 类域论已成功推广到高维域.

除少数关于更一般情形的说明外, 这些讲义将集中于局部域与数域的 Abel 扩张这一基本情形.

在引言的其余部分, 我给出 Abel 类域论的一个简短历史概览. 通篇所说数域 $K$ 的一个扩张 $L$, 是指包含在某个固定的、含 $K$ 的代数闭域中的扩张域. 例如, 若 $K$ 是包含在 $\CC$ 中的数域, 则可考虑 $K$ 在 $\CC$ 中的扩张.

\subsection*{按分裂的素理想对扩张分类}\index{fenlie@分裂(素理想)(split (prime ideal))}

回忆(ANT, 3.34) $\mathcal{O}_{K}$ 的一个素理想 $\mathfrak{p}$ 在 $K$ 的扩张 $L$ 中作(分)分解
\begin{equation*}
\mathfrak{p}\mathcal{O}_{L} = \mathfrak{P}_{1}^{e_{1}} \cdots \mathfrak{P}_{g}^{e_{g}},
\tag{2}
\end{equation*}
其中 $\mathfrak{P}_{1}, \ldots, \mathfrak{P}_{g}$ 是 $\mathcal{O}_{L}$ 中满足 $\mathfrak{P}_{i} \cap \mathcal{O}_{K} = \mathfrak{p}$ 的素理想, 且 $e_{i} \geq 1$. 此外,
\begin{equation*}
n = e_{1}f_{1} + \cdots + e_{g}f_{g},
\tag{3}
\end{equation*}
这里 $n$ 是 $L$ 在 $K$ 上的次数, $f_{i}$ 是剩余类域扩张 $\mathcal{O}_{L}/\mathfrak{P}_{i} \supset \mathcal{O}_{K}/\mathfrak{p}$ 的次数. 当某个 $e_{i} > 1$ 时, 称素理想 $\mathfrak{p}$ 在 $L$ 中分歧; 当所有 $e_{i} > 1$ 时, 称它在 $L$ 中全分歧. 当对所有 $i$ 都有 $e_{i} = 1 = f_{i}$, 从而 $\mathfrak{p} = \mathfrak{P}_{1} \cdots \mathfrak{P}_{n}$ 时, 称素理想 $\mathfrak{p}$ 在 $L$ 中分裂(或强调为完全分裂). 取 $L/K$ 的一个本原元素 $\alpha \in \mathcal{O}_{L}$. 则
\[
L = K[\alpha] \simeq K[X]/(f(X)), \quad (f\ \text{为首一不可约多项式}),
\]
且分歧的素理想整除 $f$ 的判别式(故只有有限多个)\footnote{我们有
\[
\mathrm{disc}(f)=\mathrm{disc}(\mathcal{O}_{K}[\alpha]/\mathcal{O}_{K})
=(\mathcal{O}_{L}:\mathcal{O}_{K}[\alpha])^{2}\cdot\mathrm{disc}(\mathcal{O}_{L}/\mathcal{O}_{K})
\]
(见 ANT, 2.25). 因此, 若一个素理想在 $L$ 中分歧, 则它整除 $\mathrm{disc}(f)$; 但可能存在整除 $\mathrm{disc}(f)$ 却并不在 $L$ 中分歧的素理想. 例如, 设 $L$ 是由 $\sqrt{m}$ 生成的二次域. 若 $m\equiv1\mod 4$, 则 $2$ 整除 $f$ 的判别式, 但 $2$ 在 $L$ 中并不分歧. }. 此外, 不整除 $f$ 的判别式的一个素理想在 $L$ 中分裂, 当且仅当 $f$ 模 $p$ 完全分裂. 例如, 若 $(p)$ 在 $\Q[\sqrt{m}] \simeq \Q[X]/(X^{2} - m)$ 中分歧, 则 $p$ 整除 $4m$; 而不整除 $4m$ 的素理想 $(p)$ 在 $\Q[\sqrt{m}]$ 中分裂, 当且仅当 $m$ 模 $p$ 成为平方.

当 $L/K$ 为 Galois 扩张时, $e_{1} = e_{2} = \cdots = e$ 且 $f_{1} = f_{2} = \cdots = f$, 于是这些等式简化为
\begin{equation*}
\mathfrak{p}\mathcal{O}_{L} = (\mathfrak{P}_{1} \cdots \mathfrak{P}_{g})^{e}, \quad n = efg.
\tag{4}
\end{equation*}
设 $\Spl(L/K)$ 表示 $\mathcal{O}_{K}$ 中在 $L$ 中分裂的素理想集. 十九世纪末, Frobenius 用分析方法证明了如下陈述.

\begin{theorem}[FROBENIUS]\label{thm:0.1}
当 $L/K$ 为 Galois 扩张时, 集合 $\Spl(L/K)$ 在所有 $\mathcal{O}_{K}$ 的素理想集中具有密度 $1/[L:K]$.
\end{theorem}

由此可知 $\Spl(L/K)$ 决定了 $L$: 若在 $L$ 与 $L'$ 中分裂的是同一批素理想, 则推出
\[
\Spl(L/K) = \Spl(LL'/K) = \Spl(L'/K)
\]
(ANT, 8.38 的证明), 从而
\[
[L:K] = [LL':K] = [L':K].
\]
于是 $K$ 的 Galois 扩张由集合 $\Spl(L/K)$ 分类. 仍有以下两个基本问题.

\begin{npara}\label{para:0.2} 当 $L$ 跑遍 $K$ 的有限 Galois 扩张时, 确定这些集合 $\Spl(L/K)$.
\end{npara}

\begin{npara}\label{para:0.3} 用集合 $\Spl(L/K)$ 描述扩张 $L/K$.
\end{npara}

我们将看到, 当 $L$ 跑遍 $K$ 的 Abel 扩张时, \ref{para:0.2} 有精确的答案: 此时集合 $\Spl(L/K)$ 由某些同余条件刻画. 然而对非 Abel 的 Galois 扩张, 这些集合不被同余条件刻画, 所能期望的最好结果也不过是某种解析描述.

关于 \ref{para:0.3}, 注意 $\Spl(L/K)$ 的密度决定了 $L/K$ 的次数. 我们还希望确定 $L/K$ 的整个 Galois 群, 并且希望确定 $K$ 的每个素点如何在 $L$ 中分解. 例如, 我们希望确定在 $L$ 中分歧的素理想集, 并对那些不分歧的素理想, 确定 $\mathfrak{p}$ 上方的那些素点的剩余类次数 $f(\mathfrak{p})$.

在继续之前我作一点说明: 在上述讨论中, 我们可以忽略 $K$ 的有限个素理想. 例如, 若 $S$ 是 $\mathcal{O}_{K}$ 的素理想的一个固定有限集, 则我们可以将上述讨论中的 $\Spl(L/K)$ 替换为不在 $S$ 中且在 $L$ 中分裂的素理想集 $\Spl_{S}(L/K)$.

\begin{NOTES}
\phantomsection\label{pg:orig3}对一个数域的 Galois 扩张 $L/K$, Frobenius 给 $K$ 中每个在 $L$ 中不分歧的素理想 $\mathfrak{p}$ 对应了 $\Gal(L/K)$ 中的一个共轭类 $(\mathfrak{p}, L/K)$; 他猜想, 给出固定共轭类 $C$ 的素理想的密度是 $|C|/|G|$. 按构造, $(\mathfrak{p}, L/K)$ 中的元素在 $G$ 中的阶为 $f(\mathfrak{p})$, 因此把这一陈述用于平凡共轭类即给出定理\ref{thm:0.1}. Frobenius 只能证明比其猜想更弱的一个陈述(对 \ref{thm:0.1} 已足够), 其中某些共轭类被归并到一起\footnote{Frobenius, G., 1896, \emph{Die Beziehungen zwischen den Primidealen eines algebraischen K\"orpers und den Substitutionen einer Gruppe}, Berlin Akad.-Ber. 关于 Frobenius 定理的陈述, 见\hyperlink{item:viii.A-10}{练习 A-10} (原书第 268 页).}, 而完整的猜想直到 1926 年才由 Chebotarev 证明.
\end{NOTES}

\subsection*{$\Q$ 的二次扩张}\index{ercikuozhang@二次扩张(quadratic extension)}

对 $\Q$ 的二次扩张, \ref{para:0.2} 的答案由二次互反律给出. 例如, 设 $p$ 为奇素数, 令 $p^{*} = (-1)^{\frac{p-1}{2}}p$, 于是 $p^{*} \equiv 1 \pmod{4}$. 则 $\Q[\sqrt{p^{*}}]/\Q$ 仅在 $p$ 处分歧. 一个素数 $q \neq p$ 在 $\Q[\sqrt{p^{*}}]$ 中分裂, 当且仅当 $p^{*}$ 是模 $q$ 的平方, 即 $\left(\frac{p^{*}}{q}\right) = 1$. 但当 $q$ 为奇数时, 二次互反律说 $\left(\frac{p^{*}}{q}\right) = \left(\frac{q}{p}\right)$, 故 $q$ 在 $\Q[\sqrt{p^{*}}]$ 中分裂当且仅当 $q$ 是模 $p$ 的平方. 设 $H \subset (\Z/p\Z)^{\times}$ 为平方元集. 则 $\Spl(\Q[\sqrt{p^{*}}]/\Q)$ 由满足 $q \bmod p \in H$ 的素数 $q$ 组成. 回忆 $(\Z/p\Z)^{\times}$ 是阶为 $p-1$ 的循环群, 故 $H$ 是 $(\Z/p\Z)^{\times}$ 中唯一的指数为 $2$ 的子群. 于是
\[
(\Z/p\Z)^{\times}/H \simeq \Gal(\Q[\sqrt{p^{*}}]/\Q).
\]
更一般地, 设 $S$ 为素数集的一个有限集. 令 $m$ 为只被 $S$ 中素数整除的正整数, 并设 $m$ 为奇数或能被 $4$ 整除. 对 $(\Z/m\Z)^{\times}$ 中一个指数为 $2$ 的子群 $H$, 令
\[
p(H) = \{(p) \mid p\ \text{为素数},\ p \notin S,\ p \bmod m \in H\}.
\]
则当 $L$ 跑遍在 $S$ 外不分歧的 $\Q$ 的二次扩张时, 集合 $\Spl_{S}(L/\Q)$ 恰为这些集合 $p(H)$ (对变动的 $m$). 特别地, 我们看到每个这样的集合 $\Spl_{S}(L/\Q)$ 都由某个 $m$ 模 $m$ 的同余条件决定.

\begin{NOTES}
在初等数论的课程中, 二次互反律往往显得像是一件奇特的珍玩, 其深刻程度并不超过许多其他奇特的结论. 事实上, 正如上述讨论所表明的, 它应当被视为类域论中的第一个结果, 这有助于解释 Euler, Lagrange, Legendre 与 Gauss 为何对它倾注了如此多的注意力.
\end{NOTES}

\subsection*{非分歧 Abel 扩张}\index{feifenqikuozhang@非分歧扩张(unramified extension)}\index{aboerkuozhang@Abel 扩张(abelian extension)}

设 $I$ 为 $K$ 的分式理想群, 即由 $\mathcal{O}_{K}$ 的素理想生成的自由 Abel 群; 设 $i : K^{\times} \to I$ 为把 $a \in K^{\times}$ 映到主理想 $(a)$ 的映射. $K$ 的类群 $C$ 是 $I/i(K^{\times})$. 给出一个 $C$ 的子群 $H$, 等价于给出一个包含 $i(K^{\times})$ 的 $I$ 的子群 $\tilde{H}$.

我所说 $K$ 的一个``素点'', 是指 $K$ 上非平凡赋值的等价类\footnote{英文中 “place” 比 “prime” 更常用, 我以后会改用前者; 也有人用 “spot”. }. 于是每个非零素理想对应恰好一个素点, 每个嵌入 $K \hookrightarrow \R$ 对应一个, 每对共轭的非实嵌入 $K \hookrightarrow \CC$ 对应一个. 相应的素点分别称为有限素点、实素点与复素点. $K$ 的一个元素若在对应嵌入 $K \hookrightarrow \R$ 的实素点处映到 $\R$ 中的正元素, 则称它在该实素点处为正. 若 $K$ 的一个实素点 $v$ 在 $L$ 中的每个上方素点仍为实素点, 则称 $v$ 在扩张 $L/K$ 中分裂; 否则称 $v$ 在 $L/K$ 中分歧. 例如, $\Q[\sqrt{-5}]$ 在 $\Q$ 上恰好在素点 $(2)$, $(5)$ 与 $\infty$ 处分歧.

我们说有限扩张 $L/K$ 是非分歧的, 如果它在 $K$ 的所有素点处都不分歧. 这意味着 $\mathcal{O}_{K}$ 的每个素理想在 $L$ 中都不分歧, 且 $K$ 的每个实素点仍保持为实(即 $K$ 的每个实嵌入都延拓为 $L$ 的一个实嵌入).

设 $H$ 为 $K$ 的类群 $C$ 的一个子群. 若 $\mathcal{O}_{K}$ 中在 $L$ 中分裂的素理想恰为那些落在 $\tilde{H}$ 中的素理想, 则称 $K$ 的一个有限非分歧 Abel 扩张 $L$ 是 $H$ 的类域. 下一个定理对 $K$ 的非分歧 Abel 扩张同时回答了 \ref{para:0.2} 与 \ref{para:0.3}.

\begin{theorem}[FURTW\"ANGLER 1907]\label{thm:0.4}
对 $C$ 的每个子群都存在类域; 它是唯一的, 并且 $K$ 的每个有限非分歧 Abel 扩张都作为某个 $C$ 的子群的类域出现. 若 $L$ 是 $H$ 的类域, 则 $\Gal(L/K) \approx C/H$, 并且对 $K$ 的每个素理想 $\mathfrak{p}$, $f(\mathfrak{p})$ 等于 $\mathfrak{p}$ 在群 $C/H$ 中像的阶.
\end{theorem}

对应 $K$ 的有限非分歧 Abel 扩张 $L$ 的 $C$ 的子群 $H$, 是由在 $L$ 中分裂的素理想生成的子群. 类域的唯一性由 \ref{thm:0.1} 推出.

$C$ 的平凡子群的类域称为 $K$ 的 Hilbert 类域. 它是 $K$ 在 $K$ 的所有素点(包括无穷素点)处都非分歧的最大 Abel 扩张. 在其中分裂的素理想恰为主理想, 且 $\Gal(L/K) \approx C$. 例如, $K = \Q[\sqrt{-5}]$ 的类数为 $2$, 其 Hilbert 类域为 $K[\sqrt{-1}]$ (ANT, 4.11).

\begin{NOTES}
定理\ref{thm:0.4} 由 Hilbert 于 1897 年猜想, 并由他的学生 Furtw\"angler 于 1907 年证明.
\end{NOTES}

\subsection*{分歧 Abel 扩张}\index{fenqikuozhang@分歧扩张(ramified extension)}

\phantomsection\label{pg:orig4}为把定理\ref{thm:0.4} 推广到分歧扩张, 我们需要推广理想类群的概念. 为此应先考察我们已经理解的 Abel 扩张, 即 $\Q$ 的分圆扩张. 设 $m$ 为奇数或能被 $4$ 整除的正整数, 并设 $\zeta_{m}$ 为本原 $m$ 次单位根. 回忆(ANT, 第 6 章) $\Q[\zeta_{m}]$ 是 $\Q$ 的一个扩张, 其 Galois 群为 $(\Z/m\Z)^{\times}$, 其中 $[n]$ 的作用是 $\zeta_{m} \mapsto \zeta_{m}^{n}$. 此外, 在我们关于 $m$ 的条件下, 在 $\Q[\zeta_{m}]$ 中分歧的素理想恰为满足 $p \mid m$ 的理想 $(p)$ 以及 $\infty$. 我们如何能把 $(\Z/m\Z)^{\times}$ 实现为一个理想类群? 我们可以尝试把一个理想 $(p)$ 映到 $(\Z/m\Z)^{\times}$ 中 $[p]$ 的类, 但这只有当 $(p,m) = 1$ 时才可能. 于是令 $S$ 为满足条件 $p \mid m$ 的素理想 $(p)$ 的集合, 并令 $I^{S}$ 为 $\Q$ 的、由不在 $S$ 中的素理想生成的分式理想群. $I^{S}$ 的每个元素可表示为 $(r/s)$, 其中 $r$, $s$ 为与 $m$ 互素的正整数; 我们把它映到 $(\Z/m\Z)^{\times}$ 中的 $[r][s]^{-1}$. 这给出满同态 $I^{S} \to (\Z/m\Z)^{\times}$, 其核待描述. 设 $m = \prod p^{\operatorname{ord}_{p}(m)}$, 于是
\[
\Z/m\Z \simeq \prod_{p \mid m} \Z/p^{\operatorname{ord}_{p}(m)} \quad (\text{中国剩余定理}),
\]
\[
(\Z/m\Z)^{\times} \simeq \prod_{p \mid m} (\Z/p^{\operatorname{ord}_{p}(m)})^{\times}.
\]
一个理想 $(r)$(其中 $(r,m) = 1 = (s,m)$) 映到 $1$, 当且仅当 $r$ 与 $s$ 同号, 并且对一切整除 $m$ 的 $p$, $r$ 与 $s$ 在 $(\Z/p^{\operatorname{ord}_{p}(m)})^{\times}$ 中映到同一元素. 注意在 $\infty$ 处的条件是 $\frac{r}{s} > 0$, 而在 $p$ 处的条件是 $\operatorname{ord}_{p}(\frac{r}{s} - 1) \geq \operatorname{ord}_{p}(m)$.

为推广这一点, 定义数域 $K$ 的一个\emph{模}\index{mo@模(modulus)(modulus)}\index{shexianleiqun@射线类群(ray class group)}为一个(形式)乘积 $\mathfrak{m} = \mathfrak{m}_{0}\mathfrak{m}_{\infty}$, 其中 $\mathfrak{m}_{0}$ 是 $\mathcal{O}_{K}$ 的一个整理想, $\mathfrak{m}_{\infty}$ 是 $K$ 的若干实素点的乘积. 设 $S(\mathfrak{m})$ 为整除 $\mathfrak{m}$ 的素点集, 并定义 $I^{S(\mathfrak{m})}$ 为 $\mathcal{O}_{K}$ 的、由不在 $S(\mathfrak{m})$ 中的素理想生成的分式理想群. 以\emph{模}\index{mo@模(modulus)(modulus)} $\mathfrak{m}$ 为模的\emph{射线类群}\index{shexianleiqun@射线类群(ray class group)}\index{shexianleiqun@射线类群(ray class group)} $C_{\mathfrak{m}}$ 定义为: 把 $I^{S(\mathfrak{m})}$ 模去由如下元素 $a$ 生成的主理想子群所得到的商, 其中 $a$ 满足: 在所有整除 $\mathfrak{m}_{\infty}$ 的实素点处 $a > 0$, 并且对所有整除 $\mathfrak{m}_{0}$ 的素理想 $\mathfrak{p}$ 有 $\operatorname{ord}_{\mathfrak{p}}(a - 1) \geq \operatorname{ord}_{\mathfrak{p}}(\mathfrak{m}_{0})$(这里 $\operatorname{ord}_{\mathfrak{p}}(\mathfrak{m}_{0})$ 是 $\mathfrak{p}$ 在 $\mathfrak{m}_{0}$ 中的指数).

例如, 若 $\mathfrak{m} = 1$, 则 $C_{\mathfrak{m}} = C$, 即通常的类群. 若 $\mathfrak{m}$ 是 $K$ 的实素点的乘积(于是 $\mathfrak{m}_{0} = 1$), 则 $C_{\mathfrak{m}}$ 是 $I$ 模去由 $K$ 的全正元\footnote{称 $K$ 的元素 $a$ 是\emph{全正}的, 若对 $K$ 的每个实嵌入 $\sigma\colon K\hookrightarrow\R$ 都有 $\sigma(a)>0$. 例如 $-5$ 在 $\Q[i]$ 中全正, 但在 $\Q$ 中不全正. 这看起来有些奇怪, 但就我书架上的书而言, 这个定义是标准的. 除该域为全实域的情形外, 更好的说法是: 该元素“在所有实素点处为正”. }生成的理想所得的子群的商, 即窄类群\footnote{英文原词 narrow-class group 中的连字符是要表明“窄”修饰的是那些类, 而不是这个群. }. 若 $K = \Q$ 且 $\mathfrak{m} = (m)\infty$, 则 $C_{\mathfrak{m}} \simeq (\Z/m\Z)^{\times}$, 如第一段所述.

设 $H$ 为某个模 $\mathfrak{m}$ 的射线类群 $C_{\mathfrak{m}}$ 的一个子群, 并设 $\tilde{H}$ 为它在 $I^{S(\mathfrak{m})}$ 中的原像. 称在整除 $\mathfrak{m}$ 的素点处非分歧的 $K$ 的 Abel 扩张 $L$ 是 $H$ 的类域, 若 $\mathcal{O}_{K}$ 中不整除 $\mathfrak{m}_{0}$ 且在 $L$ 中分裂的那些素理想恰为落在 $\tilde{H}$ 中的素理想. 下一个定理把定理\ref{thm:0.4} 推广到 $K$ 的所有 Abel 扩张.

\begin{theorem}[TAKAGI]\label{thm:0.5}
对射线类群 $C_{\mathfrak{m}}$ 的每个子群都存在类域; 它是唯一的, 并且 $K$ 的每个有限 Abel 扩张都作为某个射线类群的某个子群的类域出现. 若 $L$ 是 $H \subset C_{\mathfrak{m}}$ 的类域, 则 $\Gal(L/K) \approx C_{\mathfrak{m}}/H$, 并且 $K$ 的不整除 $\mathfrak{m}$ 的素理想 $\mathfrak{p}$ 在 $L$ 中不分歧, 其剩余类次数 $f(\mathfrak{p})$ 等于 $\mathfrak{p}$ 在群 $C_{\mathfrak{m}}/H$ 中像的阶.
\end{theorem}

对 $K$ 的任意有限素集 $S$, 定理\ref{thm:0.5} 完全解决了确定仅在 $S$ 中的素处分歧的 Abel Galois 扩张 $L/K$ 的集合 $\Spl_S(L/K)$ 这一问题.

该定理可以更精确. 设 $\mathfrak{m}$ 与 $\mathfrak{m}'$ 为 $K$ 的模. 若 $\mathfrak{m}' \mid \mathfrak{m}$, 则包含 $I^{S(\mathfrak{m})} \hookrightarrow I^{S(\mathfrak{m}')}$ 定义满同态 $C_{\mathfrak{m}} \twoheadrightarrow C_{\mathfrak{m}'}$; 若 $H \subset C_{\mathfrak{m}}$ 是某个子群 $H' \subset C_{\mathfrak{m}'}$ 的原像, 则 $H'$ 的每个类域也将是 $H$ 的类域. 若 $L$ 是 $H \subset C_{\mathfrak{m}}$ 的类域, 且 $H$ 并非以这种方式来自某个真整除 $\mathfrak{m}$ 的模, 则 $S(\mathfrak{m})$ 恰由在 $L$ 中分歧的素理想组成. 这个尽可能小的模称为 $L/K$ 的\emph{导子}\index{daozi@导子(conductor)}.

设 $L$ 为 $K$ 的一个 Abel 扩张. 则对每个只被在 $L$ 中分歧的素点整除的、充分大的模 $\mathfrak{m}$, $L$ 都是由在 $L$ 中分裂的素理想的类生成的子群 $H \subset C_{\mathfrak{m}}$ 的类域.

对 $K = \Q$, 该定理由 ANT 第 6 章证明的结果推出, 除了``$\Q$ 的每个 Abel 扩张都作为类域出现''这一断言, 它本质上就是 Kronecker--Weber 定理: $\Q$ 的每个 Abel 扩张都包含于某个分圆扩张.

\begin{example}\label{ex:0.6}
域 $\Q[\sqrt{6}]$ 的类数为 $1$, 故等于它自己的 Hilbert 类域. 然而它的窄类群阶为 $2$, 并且确实有 $\Q[\sqrt{6}]$ 的一个二次扩张, 即 $\Q[\sqrt{-2}, \sqrt{-3}]$, 在它所有有限素点处都非分歧于 $\Q[\sqrt{6}]$(但在两个无穷素点处都分歧).
\end{example}

当 Takagi 证明定理\ref{thm:0.5} 时, 他只证明了群 $C_{\mathfrak{m}}/H$ 与 $\Gal(L/K)$ 同构——他并未构造一个具体的同构. 有一个明显的映射 $I^{S(\mathfrak{m})} \to \Gal(L/K)$, 即把$S(\mathfrak{m})$ 中的一个素理想 $\mathfrak{p}$ 映到 $\Gal(L/K)$ 中相应的 Frobenius 元. Frobenius 定理表明这是满射, 问题是要证明它对某个 $\mathfrak{m}$ 经由 $C_{\mathfrak{m}}$ 分解, 即该映射容许一个模. Artin 于 1923 年猜想这一点, 并于 1927 年证明(见下).

\begin{example}\label{ex:0.7}
设 $d = p_{1}^{*} \cdots p_{t}^{*}$, 其中 $p_{i}$ 为互不相同的奇素数, 且 $p_{i}^{*} = (-1)^{\frac{p_{i}-1}{2}}p_{i}$ 如前. 域 $K = \Q[\sqrt{d}]$ 恰好在素理想 $(p_{1}), \ldots, (p_{t})$ 处分歧. 考虑 $K[\sqrt{p_{i}^{*}}]$. 它包含 $\Q[\sqrt{p_{1}^{*} \cdots p_{i-1}^{*} p_{i+1}^{*} \cdots p_{t}^{*}}]$, 后者在 $(p_{i})$ 处非分歧, 故 $(p_{i})$ 不能在 $K[\sqrt{p_{i}^{*}}]$ 中全分歧. 由于 $p_{i}$ 在 $K/\Q$ 中分歧, 推出 $K$ 中位于 $(p_{i})$ 之上的素点在 $K[\sqrt{p_{i}^{*}}]$ 中不分歧. 没有其他素点分歧, 故 $K[\sqrt{p_{i}^{*}}]$ 在 $K$ 上非分歧. 由 Kummer 理论(FT 5.29)知
\[
L \overset{\mathrm{def}}{=} K[\sqrt{p_{1}^{*}}, \sqrt{p_{2}^{*}}, \ldots, \sqrt{p_{t-1}^{*}}]
\]
在 $K$ 上的次数为 $2^{t-1}$, 且 $\Gal(L/K) \approx (\Z/2\Z)^{t-1}$. 设 $C_{\infty}$ 为 $K$ 的窄类群. 上面的构造表明 $(C_{\infty} : C_{\infty}^{2}) \geq 2^{t-1}$. 事实上, 只需稍多努力就能证明如下陈述:

设 $K$ 是在其中分歧的有限素点有 $t$ 个的 $\Q$ 的二次扩张. 则 $(C_{\infty} : C_{\infty}^{2}) = 2^{t-1}$ (\cite{Koch1992}, 定理 2.114).

这一结果 Gauss 已知(用不同的方法, 不同的语言). 特别地, 我们看到借助类域论, 容易构造出满足 $(C : C^{2})$ 很大的 $\Q$ 的二次扩张. 相比之下, 截至 1991 年, 尚不知道哪个二次域满足 $(C : C^{3}) > 3^{6}$. 所有在二次数域的类群中构造 $3$ 阶元素的方法似乎都涉及椭圆曲线.
\end{example}

\begin{NOTES}
射线类群 $C_{\mathfrak{m}}$ 由 Weber 于 1897 年引入. 定理\ref{thm:0.5} 由 Hilbert 与 Weber 猜想, 并由 Takagi 在 1915 至 1922 年间发表的一系列论文中证明——尤其参见他在 1920 年国际数学家大会上的报告\footnote{见他的全集. }, 其中他陈述的结果几乎与我们的完全一致.
\end{NOTES}

% ============ 以下为原书引言 pp.16--33 概述子节(此前缺失, 现补译) ============
\setcounter{theorem}{7}

\subsection*{密度定理}\index{midu@密度定理(density theorem)}\index{Dirichlet 定理@Dirichlet 定理(Dirichlet's theorem)}

设 $m$ 为 $\geq 2$ 的整数. 设 $a$ 为与 $m$ 互素的整数, 考虑序列
\begin{equation*}
\ldots, a-m, a, a+m, a+2m, a+3m, \ldots, a+km, \ldots.\tag{5}
\end{equation*}
此序列只依赖于 $a \bmod m$, 故恰有 $\varphi(m) \overset{\mathrm{def}}{=} |(\mathbb{Z}/m\mathbb{Z})^{\times}|$ 个不同的序列. Dirichlet 证明素数在这些序列中均匀分布, 即每个序列中的素数集在所有素数集中具有密度 $1/\varphi(m)$.

我们想在 Galois 群的意义下解释这一点. 设 $L/K$ 是数域的有限扩张, Galois 群为 $G$. 回忆(ANT, 第 8 章)对 $\mathcal{O}_{L}$ 中每个素理想 $\mathfrak{P}$, 存在 $G$ 中一个 Frobenius 元 $\sigma$ 使得
\begin{itemize}
\item $\sigma\mathfrak{P} = \mathfrak{P}$, 且
\item 对所有 $\alpha \in \mathcal{O}_{L}$, $\sigma\alpha \equiv \alpha^{q} \bmod \mathfrak{P}$, 其中 $q$ 是剩余类域 $\mathcal{O}_{K}/\mathfrak{p}$ 的元素个数, $\mathfrak{p} = \mathfrak{P} \cap \mathcal{O}_{K}$.
\end{itemize}
当 $\mathfrak{P}$ 在 $\mathfrak{p}$ 上不分歧时, 此元素由这些条件唯一确定——此时记作 $(\mathfrak{P}, L/K)$. 若 $\mathfrak{P}'$ 是 $\mathcal{O}_{L}$ 中另一个整除 $\mathfrak{p}$ 的素理想, 则 $\mathfrak{P}' = \tau\mathfrak{P}$ 对某个 $\tau \in G$, 且
\[
(\mathfrak{P}', L/K) = (\tau\mathfrak{P}, L/K) = \tau(\mathfrak{P}, L/K)\tau^{-1}.
\]
因此, 对 $K$ 中每个在 $L$ 中不分歧的素理想 $\mathfrak{p}$, 集合
\[
\{(\mathfrak{P}, L/K) \mid \mathfrak{P} \cap \mathcal{O}_{K} = \mathfrak{p}\}
\]
是 $G$ 中一个完整的共轭类——记作 $(\mathfrak{p}, L/K)$. 当 $G$ 交换时, 共轭类只含一个元素, 故 $(\mathfrak{p}, L/K)$ 可视为 $G$ 的元素.

现在考虑 $\mathbb{Q}[\zeta_{m}]/\mathbb{Q}$, 其中 $m$ 为奇数 $>1$ 或能被 $4$ 整除的正整数. 回忆 $\Gal(\mathbb{Q}[\zeta_{m}]/\mathbb{Q}) \simeq (\mathbb{Z}/m\mathbb{Z})^{\times}$, 其中 $[n]$ 的作用为 $\zeta \mapsto \zeta^{n}$. 不整除 $m$ 的素理想 $(p)$ 在 $\mathbb{Q}[\zeta_{m}]$ 中不分歧, 且 $(p, \mathbb{Q}[\zeta_{m}]/\mathbb{Q}) = [p]$(正如所料). 令 $\tau = [a] \in \Gal(\mathbb{Q}[\zeta_{m}]/\mathbb{Q})$. 则 $(p, \mathbb{Q}[\zeta_{m}]/\mathbb{Q}) = \tau$ 当且仅当 $p$ 落在序列 (5) 中, 于是 Dirichlet 定理说
\[
\{p \mid (p, \mathbb{Q}[\zeta_{m}]/\mathbb{Q}) = \tau\}
\]
在所有素数中具有密度 $1/\varphi(m)$. 对 $\mathbb{Q}$ 的分圆扩张, 这证明了 Frobenius 的猜想(本书 p.\pageref{pg:orig3}).

为证明他关于算术级数中的素数的定理, Dirichlet 引入了如今所谓的 Dirichlet $L$-级数. 设 $\chi \colon (\mathbb{Z}/m\mathbb{Z})^{\times} \to \mathbb{C}^{\times}$ 是 $(\mathbb{Z}/m\mathbb{Z})^{\times}$ 的一个特征(即同态). 则附于 $\chi$ 的 Dirichlet $L$-级数为
\[
L(s, \chi) = \prod_{(p,m)=1,\,p\ \text{素数}} \frac{1}{1 - \chi(p)p^{-s}} = \sum_{(n,m)=1,\,n>0} \chi(n)n^{-s}.
\]
当 $\chi$ 为平凡(主)特征时, 此式给出 Riemann $\zeta$ 函数, 只是缺了有限多个因子. Dedekind 把 $\zeta$ 函数的概念推广到每个数域:
\[
\zeta_{K}(s) = \prod_{\mathfrak{p}} \frac{1}{1 - \mathbb{N}\mathfrak{p}^{-s}} = \sum_{\mathfrak{a} \subset \mathcal{O}_{K}} \mathbb{N}\mathfrak{a}^{-s}, \quad \mathbb{N}\mathfrak{a} = (\mathcal{O}_{K} : \mathfrak{a}),
\]
而 Weber 把 $L$-级数的概念推广到任意射线类群 $C_{\mathfrak{m}}$ 的特征 $\chi$:
\[
L(s, \chi) = \prod_{(\mathfrak{p},\mathfrak{m})=\mathcal{O}_{K}} \frac{1}{1 - \chi(\mathfrak{p})\mathbb{N}\mathfrak{p}^{-s}} = \sum_{(\mathfrak{a},\mathfrak{m})=1} \chi(\mathfrak{a})\mathbb{N}\mathfrak{a}^{-s}.
\]
Weber 证明 $L(s,\chi)$ 可解析延拓到整个复平面上的亚纯函数, 并且在 $\chi$ 不是平凡特征时甚至是全纯的.

在整个复平面上这个函数仍记作 $L(s,\chi)$, 称为 Dirichlet 或 Weber $L$-函数.

\subsection*{Artin $L$-函数}\index{artinl@Artin $L$-函数(Artin $L$-function)}

设 $L/K$ 是有限 Galois 扩张, Galois 群为 $G$, 并设 $S$ 是 $\mathcal{O}_{K}$ 的有限素理想集, 包含所有在 $L$ 中分歧者. 设
\[
\rho \colon G \to \mathrm{GL}_{n}(\mathbb{C})
\]
是 $G$ 的一个表示(即群同态). 对每个不在 $S$ 中的素理想, 我们得到 $G$ 的 Frobenius 元的一个共轭类 $(\mathfrak{p}, L/\mathbb{Q})$, 它被 $\rho$ 映到 $\mathrm{GL}_{n}(\mathbb{C})$ 中的一个共轭类 $\Phi_{\mathfrak{p}}(\rho)$. $\Phi_{\mathfrak{p}}(\rho)$ 中的元素可对角化(Maschke 定理, GT 7.5), 且 $\Phi_{\mathfrak{p}}(\rho)$ 由其中元素的公共特征多项式 $\det(I - \Phi_{\mathfrak{p}}(\rho)T)$ 决定. Artin 定义
\[
L_{S}(s, \rho) = \prod_{\mathfrak{p} \notin S} \frac{1}{\det(I - \Phi_{\mathfrak{p}}(\rho)(\mathbb{N}\mathfrak{p})^{-s})}, \quad s \in \mathbb{C}.
\]
此乘积在某个右半平面中收敛到一个全纯函数. 这类函数称为 Artin $L$-函数.

注意 $G$ 的一维表示就是一个特征 $\chi \colon G \to \mathbb{C}^{\times}$, 且
\[
L_{S}(s, \chi) = \prod_{\mathfrak{p} \notin S} \frac{1}{1 - \chi(\mathfrak{p})(\mathbb{N}\mathfrak{p})^{-s}}.
\]
关于 Dirichlet 级数的一条初等等价引理表明, 因子 $\frac{1}{\det(I - \Phi_{\mathfrak{p}}(\rho)(\mathbb{N}\mathfrak{p})^{-s})}$ 由全纯函数 $L_{S}(s, \rho)$ 唯一确定. 若 $\rho$ 是单射, 则
\[
\operatorname{Spl}_{S}(L/K) = \{\mathfrak{p} \mid \Phi_{\mathfrak{p}}(\rho) = \{I\}\},
\]
故函数 $L_{S}(s, \rho)$ 决定了 $\operatorname{Spl}_{S}(L/K)$. 对非 Abel 扩张, 与其直接研究 $\operatorname{Spl}_{S}(L/K)$, 不如研究 Artin $L$-函数 $L_{S}(s, \rho)$.

\begin{NOTES}
关于 $L$-级数、尤其是 Artin $L$-级数的历史介绍, 见 Noah Snyder, \emph{Artin's $L$-functions: A Historical Approach}, 2002\cite{Snyder2002}.
\end{NOTES}

\subsection*{Artin 映射}\index{artinys@Artin 映射(Artin map)}\index{hufanying@互反映射(reciprocity map)}

回忆对 $\mathbb{Q}$ 的分圆扩张 $\mathbb{Q}[\zeta_{m}]$, 存在一个同构
\[
(\mathbb{Z}/m\mathbb{Z})^{\times} \to \Gal(\mathbb{Q}[\zeta_{m}]/\mathbb{Q}),
\]
把 $[p]$ 映到 $(p, \mathbb{Q}[\zeta_{m}]/\mathbb{Q})$. $\Gal(\mathbb{Q}[\zeta_{m}]/\mathbb{Q})$ 的一个特征 $\chi$ 通过复合给出一个 $(\mathbb{Z}/m\mathbb{Z})^{\times}$ 的特征 $\chi'$. 由于
\[
\chi'(p) = \chi((p, \mathbb{Q}[\zeta_{m}]/\mathbb{Q})),
\]
我们有
\[
L_{S}(s, \chi) = L_{S}(s, \chi').
\]
于是, 在此情形 Artin $L$-级数 $L_{S}(s, \chi)$ 是一个 Dirichlet $L$-级数.

更一般地, 类似的事是否也成立? 例如, 由一维表示定义的 Artin $L$-级数是否是一个 Weber $L$-级数? 设 $L/K$ 是 Abel 扩张, Galois 群为 $G$, 并设 $S$ 是 $K$ 的素点集, 包含那些在 $L$ 中分歧者. 设 $I^{S}$ 是由不在 $S$ 中的素理想生成的 $K$ 的分式理想群. 则 $I^{S}$ 是自由群, 故映射 $\mathfrak{p} \mapsto (\mathfrak{p}, L/K)$ 唯一地延拓为一个同态
\begin{equation*}
I^{S} \to \Gal(L/K),\tag{6}
\end{equation*}
由 Frobenius 密度定理\ref{thm:0.1}知其满射. 若存在模 $\mathfrak{m}$ 使 $S(\mathfrak{m}) = S$, 使得此同态经 $C_{\mathfrak{m}}$ 分解, 则 $\Gal(L/K)$ 的一个特征 $\chi$ 通过复合给出 $C_{\mathfrak{m}}$ 的一个特征 $\chi'$, 且
\[
L_{S}(s, \chi) = L_{S}(s, \chi'),
\]
即 Artin $L$-级数 $L_{S}(s, \chi)$ 将是一个 Weber $L$-级数.

Artin 当然猜想这样的模确实存在, 但他直到看到 Chebotarev 对其密度定理的证明后才得以证明, 并由此得到如下更完整的定理. 在陈述之前, 我们需要一些记号. 对模 $\mathfrak{m}$, 令 $P_{\mathfrak{m}}$ 为 $I^{S(\mathfrak{m})} \to C_{\mathfrak{m}}$ 的核, 于是 $P_{\mathfrak{m}}$ 由满足某些依赖于 $\mathfrak{m}$ 的正性与同余条件的元素所生成的主理想组成. 对 $K$ 的扩张 $L$, 令 $S'(\mathfrak{m})$ 表示 $L$ 中位于 $S(\mathfrak{m})$ 中素点之上的素点. 存在一个范映射 $\Nm \colon I^{S'(\mathfrak{m})} \to I^{S(\mathfrak{m})}$(ANT, p. 68), 使得 $\Nm(\mathfrak{P}) = \mathfrak{p}^{f}$, 其中 $\mathfrak{p} = \mathfrak{P} \cap \mathcal{O}_{K}$, $f$ 为剩余类次数. 注意
\[
(\Nm(\mathfrak{P}), L/K) = (\mathfrak{p}^{f}, L/K) \overset{\mathrm{def}}{=} (\mathfrak{p}, L/K)^{f} = 1,
\]
故 $\Nm(I^{S'(\mathfrak{m})})$ 包含在 $I^{S(\mathfrak{m})} \to C_{\mathfrak{m}}$ 的核中.

\begin{theorem}[ARTIN 1927]\label{thm:0.8}
设 $L$ 是 $K$ 的 Abel 扩张, $S$ 为在 $L$ 中分歧的素点集. 则对某个满足 $S(\mathfrak{m}) = S$ 的模 $\mathfrak{m}$, 同态
\[
\mathfrak{p} \mapsto (\mathfrak{p}, L/K) \colon I^{S} \to \Gal(L/K)
\]
经过 $C_{\mathfrak{m}} \overset{\mathrm{def}}{=} I^{S(\mathfrak{m})}/P_{\mathfrak{m}}$ 分解, 并定义一个同构
\[
I^{S(\mathfrak{m})}/P_{\mathfrak{m}} \cdot \Nm(I^{S'(\mathfrak{m})}) \to \Gal(L/K).
\]
特别地, 在 $L$ 中分裂的素理想恰为子群
\[
\tilde{H} = P_{\mathfrak{m}} \cdot \Nm(I^{S'(\mathfrak{m})})
\]
中的那些(属于 $I^{S}$).
\end{theorem}

定理中的同态 $\psi_{L/K} \colon C_{\mathfrak{m}} \to \Gal(L/K)$ 称为 Artin 映射(或互反映射). 正如 Artin 所指出的, 该定理包含了所有已知的互反律(见第 \ref{ch:VIII} 章), 并且应被视为这些定律对无 $1$ 次单位根之域 $K$ 的正确推广.

注意, 若 Artin 映射不是满射, 则存在 $K$ 的一个真扩张, 其中 $K$ 在 $S$ 外的每个素点都分裂. 分析方法——具体说 Frobenius 的密度定理——表明这样的扩张不存在, 但要代数地证明这一点却很困难(见第 \ref{ch:VII} 章).

由于一维表示的 Artin $L$-函数都是 Weber $L$-级数, 它们可延拓为整个复平面上的亚纯函数. Artin 猜想有限群的每个特征都是接近 Abel 的特征的 $\mathbb{Z}$-线性组合, 并证明此猜想蕴含所有 Artin $L$-级数都延拓为整个复平面上的亚纯函数. 这一点由 Brauer 于 1946 年证明.

Artin 猜想, 若 $\rho$ 是非平凡不可约表示, 则 $L(s, \rho)$ 在整个复平面上全纯. 他指出证明此猜想需要全新的方法. 这一现在称为 Artin 猜想的命题, 是数论中最重要的猜想之一. 在定理 1.8 之后, 一维表示的 Artin $L$-函数是 Weber $L$-函数, 故 Artin 猜想由 Weber 的结果推出. 一般地, 如 Langlands 所解释的, 不可约表示的 Artin $L$-函数应是尖形式自守 $L$-函数, 而后者已知是全纯的. 对二维表示, Artin 猜想大体上已被证明(Artin, Langlands, Tunnell, Taylor 等). 参见 Taylor, Richard, \emph{Pacific J. Math.} 181, 1997, 337–347, 及其后续论文.

对函数域, Weil 在证明有限域上曲线的 Riemann 假设的同时证明了 Artin 猜想, 并认为这是更重要的结果.

\subsection*{局部类域论与无限扩张}\index{juubu@局部类域论(local class field theory)}\index{wuxiankuozhang@无限扩张(infinite extension)}

$K$ 的 Abel 扩张 $L$ 定义 $K$ 在素点 $v$ 处的完备化 $K_{v}$ 的一个 Abel 扩张 $L^{v}/K_{v}$. 反之, $K_{v}$ 的每个有限扩张都来自 $K$ 的同一个次数的扩张(ANT, 7.62). 因此, 数域的 Abel 扩张的分类中包含局部域的 Abel 扩张的分类, 这并不令人意外. Hasse (1930)\footnote{H. Hasse, \emph{Die Normenresttheorie relative-Abelscher Zahlkörper als Klassenkörper im Kleinen}, J. für Mathematik (Crelle) 162 (1930), 145–154.} 明确表述了这一点.

\begin{theorem}\label{thm:0.9}
设 $K$ 是 $\mathbb{Q}_{p}$ 的有限扩张. 对 $K$ 的每个有限 Abel 扩张 $L$, $\Nm(L^{\times})$ 是 $K^{\times}$ 中有限指数的子群, 且映射 $L \mapsto \Nm(L^{\times})$ 是从 $K$ 的有限 Abel 扩张集到 $K^{\times}$ 的有限指数子群集上的一个反序双射.
\end{theorem}

于是到 1930 年, 数域与局部域的 Abel 扩张都已被分类. 然而, 该理论有三方面被认为不够令人满意.
\begin{itemize}
\item 局部域比整体域简单, 人们期望先局部地证明一个命题, 再整体地推出, 而非相反. 至少, 应当有一个纯局部的局部类域论证明.
\item 类域论的主要陈述是代数性的, 故它们应容许一个纯代数证明.
\item 定理\ref{thm:0.8} 分类了导子整除某个固定模的、数域 $K$ 的 Abel 扩张. 要包含其他 Abel 扩张, 就不得不不断放大导子. 这就提出了如下问题: 同时给出所有 Abel 扩张的一个统一描述, 或者说, 无限 Abel 扩张的类域论.
\end{itemize}

Chevalley 解决了所有这三个问题. 令 $K^{\mathrm{ab}}$ 表示 $K$ 的所有有限 Abel 扩张的复合. 局部类域论的完满陈述是: 对每个 $p$-进域 $K$, 存在定义好的同态 $\phi \colon K^{\times} \to \Gal(K^{\mathrm{ab}}/K)$(现称局部 Artin 映射), 它对每个有限 Abel 扩张 $L$ 诱导同构
\[
K^{\times}/\Nm(L^{\times}) \to \Gal(L/K),
\]
并且所有(开)有限指数子群都是范群. 这表明有可能定义一个整体 Artin 映射, 其各局部分量为局部 Artin 映射, 即如下图表对所有 $K$ 的素点 $v$ 交换:
\[
\begin{tikzcd}
K_{v}^{\times}\ar[r,"\phi_{v}"]\ar[d] & \Gal(K_{v}^{\mathrm{ab}}/K_{v})\ar[d]\\
\prod_{v} K_{v}^{\times}\ar[r] & \Gal(K^{\mathrm{ab}}/K).
\end{tikzcd}
\]
这个问题在于 $\prod_{v} K_{v}^{\times}$ 不是局部紧的. 紧群的乘积是紧的, 但局部紧群不满足类似结论, 而群 $K_{v}^{\times}$ 只是局部紧的. 事实上 $\prod_{v} K_{v}^{\times}$ 太大, 以致不存在这样的 $\phi$. Chevalley 通过定义 $K$ 的\emph{伊代}群 $\mathbb{I}_{K}$ 解决了这个问题: $\mathbb{I}_{K}$ 是 $\prod_{v} K_{v}^{\times}$ 中由满足对几乎所有非阿基米德 素点有 $a_{v} \in \mathcal{O}_{v}^{\times}$ 的族 $(a_{v})$ 组成的子群. 赋予使 $\prod_{v\mid\infty} K_{v}^{\times} \times \prod_{v\ \text{有限}} \mathcal{O}_{v}^{\times}$ 为开子群之拓扑时, $\mathbb{I}_{K}$ 成为局部紧群. 把 $K^{\times}$ 作为对角线子群嵌入 $\mathbb{I}_{K}$.

在 Chevalley 的类域论进路中, 先直接证明局部类域论. 然后定义一个整体 Artin 映射 $\phi_{K} \colon \mathbb{I}_{K} \to \Gal(K^{\mathrm{ab}}/K)$, 其各局部分量为局部 Artin 映射, 并证明 $K^{\times}$ 包含在 $\phi_{K}$ 的核中, 且 $\phi$ 诱导的同态
\[
\mathbb{I}_{K}/K^{\times} \to \Gal(K^{\mathrm{ab}}/K)
\]
满射, 其核等于 $\mathbb{I}_{K}/i(K^{\times})$ 的恒等连通分支. 换言之, 同态 $\phi_{K}$ 适合如下交换图表
\[
\begin{tikzcd}
K_{v}^{\times}\ar[r,"\phi_{v}"]\ar[d] & \Gal(K_{v}^{\mathrm{ab}}/K_{v})\ar[d]\\
\mathbb{I}_{K}\ar[r,"\phi_{K}"] & \Gal(K^{\mathrm{ab}}/K)
\end{tikzcd}
\]
对所有 $K$ 的素点 $v$ 成立, 且 $\phi_{K}$ 诱导同构
\[
\mathbf{C}_{K}/\mathbf{C}_{K}^{\circ} \to \Gal(K^{\mathrm{ab}}/K), \quad \mathbf{C}_{K} \overset{\mathrm{def}}{=} \mathbb{I}_{K}/K^{\times}.
\]
此陈述与定理\ref{thm:0.9} 的关系如下: 存在典范同构 $C/C^{\circ} \simeq \varprojlim_{\mathfrak{m}} C_{\mathfrak{m}}$, 并且对 $K$ 的每个有限 Abel 扩张 $L$ 与充分大的模 $\mathfrak{m}$, 有交换图表
\[
\begin{tikzcd}
\mathbf{C}\ar[r,"\phi_{K}"]\ar[d] & \Gal(K^{\mathrm{ab}}/K)\ar[d,"\tau \mapsto \tau|_{L}"]\\
C_{\mathfrak{m}}\ar[r,"\psi_{L/K}"] & \Gal(L/K).
\end{tikzcd}
\]
到 1950 年, 局部 Artin 映射已能在局部被刻画, 或被视为整体 Artin 映射的局部分支, 但直到 1965 年 Lubin 与 Tate 才找到它的一个显式局部描述.

\begin{NOTES}
在数域情形, Artin 映射 $C \to \Gal(K^{\mathrm{ab}}/K)$ 满射, 其核为 $C$ 的恒等分支. 在函数域情形, 它是单射但不满射. Weil 曾猜想, 把数域的整个伊代类群解释为 Galois 群, 可能是证明 Riemann 假设的关键(见本书 p.\pageref{pg:orig180}).
\end{NOTES}

\subsection*{中心单代数}\index{zhongxin@中心单代数(central simple algebra)}\index{Brauer 群@Brauer 群(Brauer group)}

域 $K$ 上的中心代数是指中心为 $K$ 且在 $K$ 上有限维的代数 $A$; 若 $A$ 没有非零真双侧理想, 则称其为单的. 例如, 矩阵代数 $M_{n}(K)$ 是 $K$ 上的中心单代数. 第一个有趣的中心单代数例子, 即 $\mathbb{R}$ 上的四元数代数, 由 Hamilton 于 1843 年发现. 它在 $\mathbb{R}$ 上由 $1, i, j, ij$ 张成, 其乘法由
\[
i^{2} = -1, \quad j^{2} = -1, \quad ij = -ji.
\]
决定.

设 $K$ 为域. 1906 年, Dickson 说明如何把 $K$ 的一个循环扩张\footnote{即, 具有循环 Galois 群的 Galois 扩张.} $L/K$、$\Gal(L/K)$ 的一个生成元 $\sigma$ 以及 $K$ 的一个非零元 $a$ 附到一个中心单代数 $(L/K, \sigma, a)$ 上. 作为 $L$-向量空间, 该代数由元素 $1, \alpha, \ldots, \alpha^{n-1}$ 张成, 其中 $n = [L:K]$, 其乘法由
\[
\alpha^{n} = a, \quad \alpha c = (\sigma c)\alpha \quad (\text{对 } c \in L)
\]
决定.

于是 $(L/K, \sigma, a)$ 是一个包含 $L$ 作为极大子域的中心单代数, 其中 $\sigma$ 在 $L$ 上的作用由某个内自同构诱导.

Brauer 与 Noether 说明了如何用因子系描述中心单代数, 一旦定义了群上同调, 这些因子系即被认出是 2-上循环.

1907 年, Wedderburn 证明 $k$ 上的每个中心单代数都同构于某个中心可除代数上的矩阵代数. Brauer 定义 $K$ 上的两个中心单代数相似, 如果它们同构于同一个可除代数上的矩阵代数, 并证明相似类在张量积下构成一个群(现称 $K$ 的 Brauer 群). 于是, 分类域 $K$ 上的中心单代数的问题变成了确定该域的 Brauer 群 $\mathrm{Br}(K)$.

Frobenius 的结果确定了 $\mathbb{R}$ 的 Brauer 群, Hasse 确定了非阿基米德局部域的 Brauer 群. 1932 年, Albert, Brauer, Hasse 与 Noether 证明数域上的每个可除代数都是循环的, 并确定了这样的域的 Brauer 群. 这是一个与类域论基本定理同样深刻的定理, 在这些讲义中将与它们同时证明.

Brauer 群有许多应用, 例如对有限群的表示理论(这正是 Brauer 所感兴趣的), 以及对数域上半单代数群的分类.

\begin{NOTES}
数域的 Brauer 群的确定过程在 Fenster 与 Schwermer 2005\cite{FensterSchwermer2005} 以及 Roquette 2005\cite{Roquette2005} 中有精彩的叙述. 我引用前者的 MR 评论:

In February, 1931, Hasse began a joint effort with Richard Brauer and Emmy Noether to classify all division algebras over algebraic number fields. They achieved this in a paper that was published in early 1932. The young A. A. Albert was also working on the problem at the same time. He had been in correspondence with Hasse during 1930–1931 and he had come close to resolving the problem. Some of his remarks to Hasse had been helpful in the Brauer-Hasse-Noether solution. Within weeks of that event, Albert found an alternate proof, which utilized a suggestion from Hasse.
\end{NOTES}

\subsection*{上同调}\index{shangtong@上同调(cohomology)}

对群 $G$ 与 $G$-模 $M$, 有同调群 $H_{r}(G, M)$($r \geq 0$)与上同调群 $H^{r}(G, M)$($r \geq 0$). $G$-模的短正合列给出同调与上同调群的长正合列. Tate 证明, 当 $G$ 有限时, 这些长正合列可以拼接成一条很长的序列. 更确切地说, 对所有 $r \in \mathbb{Z}$ 定义的群
\[
H_{T}^{r}(G, M) \overset{\mathrm{def}}{=}
\begin{cases}
H^{r}(G, M), & r > 0, \\
M^{G}/\Nm(M), & r = 0, \\
\operatorname{Ker}(\Nm)/I_{G}M, & r = -1, \\
H_{-r-1}(G, M), & r < -1,
\end{cases}
\]
具有此性质. 特别地,
\[
H_{T}^{-2}(G, \mathbb{Z}) \overset{\mathrm{def}}{=} H_{1}(G, \mathbb{Z}) \simeq G^{\mathrm{ab}} \quad (\text{为 } G \text{ 的最大 Abel 商}).
\]
设 $C$ 是一个 $G$-模, 使得对所有 $G$ 的子群 $H$ 都有 $H^{2}(H, C)$ 是阶为 $|H|$ 的循环群. 若 $H^{1}(G, C) = 0$, 则 Tate 证明选取 $H^{2}(G, C)$ 的一个生成元确定了同构族
\[
H_{T}^{r}(G, M) \to H_{T}^{r+2}(G, M \otimes C), \quad \text{所有 } r \in \mathbb{Z}.
\]
对局部域的 Galois 扩张 $L/K$(Galois 群为 $G$), 计算 $K$ 的 Brauer 群表明 $H^{2}(G, L^{\times})$ 是阶为 $|G|$ 的循环群, 并且有一个典范生成元. 由于 Hilbert 定理 90 给出 $H^{1}(G, L^{\times}) = 0$, Tate 定理给了我们典范同构
\[
H_{T}^{r}(G, \mathbb{Z}) \to H_{T}^{r+2}(G, L^{\times})
\]
对所有 $r \in \mathbb{Z}$. 当 $r = -2$ 时, 这成为同构
\[
G^{\mathrm{ab}} \xrightarrow{\simeq} K^{\times}/\Nm(L^{\times}),
\]
其逆就是局部 Artin 映射. 对整体域有类似的论证, 此时以伊代模去主伊代群代替域的乘法群.

\begin{NOTES}
我引用 Mac Lane 1978\cite{MacLane1978}, p. 17:

\begin{quote}
Mac Lane 回忆说, Artin(约 1948 年)在谈话中指出, 群的上同调本应在类域论中发挥作用. Hochschild (1950) 与 Hochschild 和 Nakayama (1952) 说明了类域论的 Brauer 群论证如何能被上同调论证取代. 1952 年, Tate 证明有限群 $G$ 的同调群与上同调群可以恰当地合并为一条长正合列. 他利用此序列, 连同转移与限制的性质, 给出类域论的一个优美重构.
\end{quote}
\end{NOTES}

\subsection*{类域的显式构造(Hilbert 第十二问题)}\index{hilbert12@Hilbert 第十二问题(Hilbert's 12th problem)}\index{xianshi@显式类域论(explicit class field theory)}

Kronecker--Weber 定理表明 $\mathbb{Q}$ 的每个 Abel 扩张都包含于由指数函数的一个特殊值 $\zeta = e^{2\pi i/m}$ 生成的域中. 稍后的一个定理(Kronecker 青春之梦)表明, 虚二次数域的每个 Abel 扩张都包含于由椭圆模函数的某些特殊值生成的域中. 在著名的 23 个问题中的第 12 个里, Hilbert 问: 其他数域的 Abel 扩张能否由显式全纯函数的特殊值生成.

推广椭圆模函数的函数已被充分了解, 其特殊值的算术性质亦然——这就是复乘理论与 Shimura 簇理论——但它们只对构造 CM 域的 Abel 扩张有用, 并且即使对这些域, 它们也只给出(粗略地说)不从全实子域的 Abel 扩张来的那些 Abel 扩张(除 $\mathbb{Q}$ 外). 另一方面, Stark 的猜想为全实域的 Abel 扩张的生成元提供了候选.

\subsection*{显式类域论}

类域论中定理的证明优美但抽象. 例如, 局部类域论最初的局部证明只表明存在唯一的 Artin 映射 $\phi \colon K^{\times} \to \Gal(K^{\mathrm{ab}}/K)$ 具有某些性质, 却未解决该映射与 $K^{\mathrm{ab}}$ 的显式描述问题. 所幸, 在这种情况下, 我们将在第 \ref{ch:I} 章解释的 Lubin 与 Tate 的理论对此问题给出了优美的回答.

已有大量工作致力于寻找类域论中其他映射与对象的显式描述. 在当前版本的讲义中, 这一主题大体被忽略.

\subsection*{计算类域论}\index{jisuan@计算类域论(computational class field theory)}

除了寻找类域论中映射与对象的显式描述外, 还可以要求计算它们的算法. 对此主题, 我只能指引读者参阅 Cohen 2000\cite{Cohen2000}.

\subsection*{非 Abel 类域论}\index{fiabel@非 Abel 类域论(nonabelian class field theory)}

在 1920 年 ICM 上他向世人宣布 Abel 类域论基本定理证明的同一报告中, Takagi 提出了非 Abel 类域论的问题:
\begin{quote}
En m'arrêtant ici, je me permets d'attirer votre attention sur un problème important de la théorie des nombres algébriques: à savoir, rechercher s'il est possible de définir la classe d'idéaux d'un corps algébrique de telle manière que les corps supérieur relativement normal mais non abélien puisse être caractérisé par le groupe correspondant de classes d'idéaux du corps de fond.\footnote{意即: 在此停笔, 请允许我提请诸位注意代数数论中的一个重要问题, 即能否以这样的方式定义代数数域的理想类, 使得相对正规但非 Abel 的扩张域可由基域的相应理想类群刻画.}
\end{quote}
范限制定理表明, (射线)类群的子群不能区分一个扩张域与其最大的 Abel 子扩张.

几十年来, 非 Abel 类域论应采取何种形式, 甚至它是否存在, 都不清楚. 1946 年, Artin 猜测找到正确的陈述是唯一的问题: 一旦我们知道它们是什么, 就能由 Abel 类域论推出它们(\emph{A Century of Mathematics in America}, Part II, (Peter Duren 编), 1989, p. 312). Weil 转述说, 一年后 Artin 表示他已对存在非 Abel 类域论失去信心(Weil, A., \emph{Œuvres}, Vol. III, p. 457.)

与其研究集合 $\operatorname{Spl}_{S}(L/K)$, 不如研究 $\Gal(L/K)$ 的表示 $\rho$ 的 Artin $L$-级数 $L(s, \rho)$. 描述集合 $\operatorname{Spl}_{S}(L/\mathbb{Q})$ 的问题于是变成描述以这种方式出现的解析函数集的问题. Langlands 构造了一类 $L$-级数, 称为自守 $L$-级数, 并猜想\footnote{注意它与谷山猜想的相似之处——事实上二者都是一个更广泛猜想的特例.} 每个 $L_{S}(s, \rho)$ 都是自守的, 并且指明了哪些自守 $L$-级数以这种方式出现. 于是, 此猜想对所有有限 Galois 扩张 $L/\mathbb{Q}$ 回答了原问题. 对 $n = 1$(即 $G$ Abel)与所有 $K$, Artin 证明所有 Artin $L$-级数都是自守的——这是他证明定理\ref{thm:0.8} 的动机. 如前所述, 对 $n = 2$, 该猜想在大多数情形已被证明.

1967 年, Langlands 陈述了他猜想的\footnote{见 1967 年致 Weil 的一封信, 以及后来(对我们其余人)的 \emph{Problems in the theory of automorphic forms} (1970). 对这些著作的一个引人入胜的介绍, 见 Casselman 2001\cite{Casselman2001}. Langlands 著作(包括上述两篇)的 TeX 版本可见于 \url{http://publications.ias.edu/rpl/}.} 函子性原理, 它把非 Abel 类域论作为一个特例包含在内. 对局部域 $K$, 可表述如下. $K$ 的 Weil 群 $W_{K}$ 定义为 $\Gal(K^{\mathrm{al}}/K)$ 中那些在剩余类域上作用为 Frobenius 元的整数次幂的元素组成的子群. Abel 局部类域论中的局部 Artin 映射可视为从 $K^{\times}$ 到 $W_{K}$ 的最大 Abel 商 $W_{K}^{\mathrm{ab}}$ 的同构 $\phi_{K}$. Langlands 猜想从 $W_{K}$ 到 $\mathrm{GL}_{n}(\mathbb{C})$ 的同态对应于 $\mathrm{GL}_{n}(K)$ 的某些表示. 对 $n = 1$, $\mathrm{GL}_{1}(K) = K^{\times}$ 的表示恰为特征, 并且对应由与 $\phi_{K}$ 的复合给出. 对 $n > 1$, $\mathrm{GL}_{n}(K)$ 的表示通常是无限维的.

在自守一侧, 设 $\mathcal{A}_{n}(K)$ 为 $\mathrm{GL}_{n}(K)$ 在复向量空间上的、每个向量的稳定子开的不可约表示的等价类集. 在 Galois 侧, 设 $\mathcal{G}_{n}(K)$ 为对 $(r, N)$ 的等价类集, 其中 $r$ 是 $W_{K}$ 在 $n$ 维复向量空间 $V$ 上的半单表示, 在某个开子群上平凡, $N$ 是 $V$ 的一个幂零自同态, 使得用 $r(\sigma)$($\sigma \in W_{K}$)共轭 $N$ 会乘上 $\phi_{K}^{-1}(\sigma)$ 的绝对值. $K$ 的局部 Langlands 猜想断言存在双射族 $(\sigma_{n})_{n \geq 1}$,
\[
\pi \mapsto \sigma_{n}(\pi) \colon \mathcal{A}_{n}(K) \to \mathcal{G}_{n}(K),
\]
使得
\begin{enumerate}[(a)]
\item $\sigma_{n}(\pi)$ 的行列式视为 $W_{K}$ 的特征时, 经 $\phi_{K}$ 对应于 $\pi$ 的中心特征;
\item 映射 $\sigma_{n}$ 保持各对表示的 $L$-因子与 $\varepsilon$-因子(自守侧由 Jacquet, Piatetskii-Shapiro 与 Shalika 定义, Galois 侧由 Langlands 与 Deligne 定义);
\item 对 $\chi \in \mathcal{A}_{1}(K)$, $\sigma_{n}(\pi \otimes (\chi \circ \det)) = \sigma_{n}(\pi) \otimes \sigma_{1}(\chi)$;
\item $\sigma_{n}$ 与取对偶 $\pi \mapsto \pi^{\vee}$ 交换.
\end{enumerate}
对每个 $K$, Henniart 证明这样的族至多存在一个. 该猜想本身由 Harris 与 Taylor (2001)\cite{HarrisTaylor2001} 证明. 几个月后, Henniart (2000)\cite{Henniart2000} 找到了一个更简单的证明.

\subsection*{最后的注记}\index{zhuihou@最后的注记(final remark)}

回忆 Dedekind 的如下定理(ANT, 3.41).

\begin{theorem}\label{thm:0.10}
设 $A$ 是具有分式域 $K$ 的 Dedekind 整环, $B$ 为 $K$ 在有限可分扩张 $L$ 中的整闭包. 设 $L = K[\alpha]$, $\alpha \in B$, 并设 $f(X)$ 为 $\alpha$ 在 $K$ 上的极小多项式. 则下列关于 $A$ 的素理想 $\mathfrak{p}$ 的条件等价:
\begin{enumerate}[(a)]
\item $\mathfrak{p}$ 不整除 $\operatorname{disc}(f(X))$;
\item $\mathfrak{p}$ 在 $B$ 中不分歧, 且 $A_{\mathfrak{p}}[\alpha]$ 是 $A_{\mathfrak{p}}$ 在 $L$ 中的整闭包(这里 $A_{\mathfrak{p}} = \{a/b \mid b \notin \mathfrak{p}\}$);
\item 存在分解
\[
f(X) \equiv f_{1}(X) \cdots f_{g}(X) \mod \mathfrak{p}
\]
其中诸 $f_{i}$ 互异、首一且在模 $\mathfrak{p}$ 下不可约.
\end{enumerate}
\end{theorem}

当这些条件成立时, $\mathfrak{p}$ 在 $L$ 中分解为素理想的乘积为
\[
\mathfrak{p}B = (\mathfrak{p}, f_{1}(\alpha)) \cdots (\mathfrak{p}, f_{g}(\alpha)).
\]
于是类域论实际上关于系数在数域中的一元多项式及其根: $K$ 的 Abel 扩张对应于首一不可约多项式 $f(X) \in K[X]$, 使得给出域自同构的 $f(X)$ 的根的置换可交换; Dedekind 定理表明, 在 Abel 扩张 $L$ 中, $K$ 的除了有限多个外的所有素理想的分解对应于有限域上多项式的分解.

\subsection*{本讲义采用的方法}\index{benjiangyi@本讲义采用的方法(approach taken in these notes)}

在这些讲义中, 我们遵循 Chevalley (1940, 1954)\cite{Chevalley1940}\cite{Chevalley1954}、Artin 与 Tate (1961)\cite{ArtinTate1961} 以及 Tate (1952, 1967)\cite{Tate1952}\cite{Tate1967} 的代数/上同调进路, 证明局部与整体类域论的基本定理. 上同调进路清楚地表明了类域论基本定理与 Brauer 群计算之间的关系. 事后看来, 自 Takagi 的论文以来, 上同调计算就隐含在类域论之中.

然而, 类域论还有其他进路, 当能提供额外洞见时, 我不惜给出多于一种证明. 例如, 我说明如何借助分析简化整体情形基本定理证明的某些步骤. 此外, 我说明(遵循 Lubin 与 Tate)形式群律理论如何能使局部类域论的结果更显式、更令人满意.

尽管可以既不用分析也不用上同调来证明类域论的基本定理, 但有些主要定理甚至离不开其中之一才能陈述, 例如关于素数密度的定理, 或关于数域的上同调群的定理. 近些年来, 上同调结果在类域论的许多应用中至关重要.\footnote{例如, 据 Google Scholar, 我的《算术对偶定理》一书已被 1014 篇研究论文引用(自其开始统计以来). MR 列出 430 篇(约自 2000 年起).}

课程的核心是奇数编号的章节. 第 \ref{ch:II} 章(论群的上同调)是其余各章的基础, 但第 \ref{ch:IV} 章、第 \ref{ch:VI} 章与第 \ref{ch:VIII} 章并非阅读第 \ref{ch:III}、\ref{ch:V}、\ref{ch:VII} 章所必需. 除第一节外, 第 \ref{ch:I} 章可被那些希望尽快到达整体类域论基本定理的人跳过.

\begin{exercise}\label{ex:0.11}
完成证明: 二次互反律允许描述 $L/\mathbb{Q}$ 为二次扩张时的集合 $\operatorname{Spl}_{S}(L/\mathbb{Q})$.
\end{exercise}

\begin{exercise}\label{ex:0.12}
证明 $\mathbb{Q}[\sqrt{-5}, \sqrt{-1}]$ 是 $\mathbb{Q}[\sqrt{-5}]$ 的 Hilbert 类域.
\end{exercise}

\begin{exercise}\label{ex:0.13}
证明 本书 p.\pageref{pg:orig4} 上定义的映射 $I^{S} \to (\mathbb{Z}/m\mathbb{Z})^{\times}$ 具有所描述的核, 从而诱导同构 $C_{\mathfrak{m}} \to (\mathbb{Z}/m\mathbb{Z})^{\times}$.
\end{exercise}

\begin{exercise}\label{ex:0.14}
证明例 \ref{ex:0.6} 中的陈述.
\end{exercise}

\begin{exercise}\label{ex:0.15}
设 $L = \mathbb{Q}[\sqrt{-1}, \sqrt{-5}]$. 则 $\Gal(L/\mathbb{Q}) = \{1, \sigma, \tau, \sigma\tau\} \simeq \mathbb{Z}/2\mathbb{Z} \times \mathbb{Z}/2\mathbb{Z}$, 其中 $\sigma$ 固定 $\mathbb{Q}[\sqrt{-1}]$, $\tau$ 固定 $\mathbb{Q}[\sqrt{-5}]$, $\sigma\tau$ 固定 $\mathbb{Q}[\sqrt{5}]$.
\begin{enumerate}[(a)]
\item 证明在 $L$ 中仅 $2, 5, \infty$ 分歧.
\item 对所有 $p \neq 2, 5$, 计算 $(p, L/\mathbb{Q})$.
\item 设 $\mathfrak{m} = (20)\infty$. 证明对某个 $H \subset C_{\mathfrak{m}}$, 映射 $p \mapsto (p, L/\mathbb{Q})$ 定义同构 $C_{\mathfrak{m}}/H \to \Gal(L/\mathbb{Q})$, 并求 $H$.
\end{enumerate}
提示: 证明 $L \subset \mathbb{Q}[\zeta]$, 其中 $\zeta$ 为本原 20 次单位根.
\end{exercise}

我将给你讲一个关于互反律的故事. 在我完成博士论文后, 我有了为非 Abel 扩张定义 $L$-级数的想法. 但要使它们与 Abel 扩张的 $L$-级数一致, 某个同构必须成立. 我能证明它蕴含所有标准的互反律. 于是我称之为一般互反律, 并试图证明它, 却做不到, 即便尝试多次. 然后我把它拿给其他数论学家看, 但他们都嘲笑它, 我记得尤其是 Hasse 对我说这绝不可能是真的. 尽管如此, 我仍坚持, 但毫无进展. 连续三年, 没有一周我不去尝试证明互反律. 这令人气馁, 其间我也转向其他事情. 后来有一天下午我没什么特别的事做, 于是说: “好吧, 我再试一次证明互反律.” 于是我走出去, 在花园里坐下. 你看, 从一开始我就有用分圆域的想法, 但它们从未奏效, 而此刻我突然明白, 一直以来我用错了它们——半小时内我就证出来了.

Emil Artin, 由 Mattuck 回忆(见 \emph{Recountings: Conversations with MIT Mathematicians} 2009).

为何其他数学家如此难以相信它? 当时的互反律与幂剩余符号紧密相连. 实际上 Euler 已以正确的方式表述了二次互反律(即符号 $(\Delta/p)$ 只依赖于 $p$ 模 $\Delta$ 的剩余类), 但 Legendre 的表述占了上风. 我很确信, 如果 Artin 立刻向 Hasse 表明一般互反律等价于已知的幂剩余互反律, Hasse 就不会笑. Hasse 不必等 Chebotarev 的密度定律来说服自己 Artin 是对的——到 Chebotarev 的想法出现时, Artin 必已无疑. 顺便说, Chebotarev 的文章做的远不止证明 Frobenius 元的重要性——它为 Artin 提供了证明其互反律的关键思想, 即 Hilbert 的 Abel 交叉技术.

让我再补充一点: Felix Bernstein 于 1904 年猜想了一个互反律, 在非分歧 Abel 扩张的特殊情形大体等价于 Artin 的定律. 它的技术性表明, 从已知的幂剩余互反律猜出像 Artin 那样简单的定律绝非易事.\footnote{见 Franz Lemmermeyer 在 \mo{243590} 上的讨论.}

% 恢复正文章节定理编号(供后续各章使用)
\renewcommand{\thetheorem}{\thesection.\arabic{theorem}}
